\documentclass{article}
\usepackage{amsmath,amssymb,amsthm}
\usepackage{algorithm,algorithmic}
\usepackage{hyperref}
\newtheorem{theorem}{Theorem}
\newtheorem{lemma}{Lemma}
\newtheorem{assumption}{Assumption}
\newtheorem{corollary}{Corollary}
\newcommand{\prox}{\operatorname{prox}}
\newcommand{\R}{\mathbb{R}}
\begin{document}

\section*{Proximal Gradient Descent (PGD)}

\section*{Mathematical Formulation}
Let  
\[
F(x)=f(x)+g(x),\qquad f:\R^{n}\to\R\ \text{convex and }L\text{-smooth},\;
g:\R^{n}\to\R\cup\{+\infty\}\ \text{convex, proper, l.s.c.}
\]
Given a stepsize $\eta>0$, the \emph{proximal gradient} iteration is  
\[
\boxed{%
x_{k+1}= \prox_{\eta g}\!\bigl(x_{k}-\eta\nabla f(x_{k})\bigr)},\qquad k=0,1,\dots
\tag{PGD}
\]
where the proximal operator of $g$ is  
\[
\prox_{\eta g}(v)=\arg\min_{u\in\R^{n}}\Bigl\{g(u)+\frac{1}{2\eta}\|u-v\|^{2}\Bigr\}.
\]

\section*{Pseudocode}
\begin{algorithm}[H]
\caption{Proximal Gradient Descent}
\begin{algorithmic}[1]
\STATE \textbf{Input:} $x_{0}\in\R^{n}$, stepsize $\eta\in(0,1/L]$, maximum iterations $K$.
\FOR{$k=0,\dots,K-1$}
    \STATE Compute gradient $g_{k} \leftarrow \nabla f(x_{k})$.
    \STATE Gradient step $y_{k}\leftarrow x_{k}-\eta g_{k}$.
    \STATE Proximal step $x_{k+1}\leftarrow \prox_{\eta g}(y_{k})$.
\ENDFOR
\STATE \textbf{Output:} $x_{K}$ (or the average $\bar x_{K}=\frac1K\sum_{k=1}^{K}x_{k}$).
\end{algorithmic}
\end{algorithm}

\section*{Dry Run Example}
Consider the scalar problem  
\[
f(w)=(w-3)^{2},\qquad g\equiv0,\qquad \eta=0.1,\qquad w_{0}=0 .
\]
Since $g\equiv0$, $\prox_{\eta g}$ is the identity map.  
The gradient $\nabla f(w)=2(w-3)$.  Table of three iterations:

\begin{center}
\begin{tabular}{c|c|c|c}
$k$ & $w_{k}$ & $\nabla f(w_{k})$ & $f(w_{k})$\\ \hline
0 & $0$ & $-6$ & $(0-3)^{2}=9$\\
1 & $0.6$ & $-4.8$ & $(0.6-3)^{2}=5.76$\\
2 & $1.08$ & $-3.84$ & $(1.08-3)^{2}=3.6864$\\
3 & $1.464$ & $-3.072$ & $(1.464-3)^{2}=2.3593$\\
\end{tabular}
\end{center}
The update rule $w_{k+1}=w_{k}-\eta\nabla f(w_{k})$ yields the numbers shown, and the objective value decreases monotonically.

\newpage

\section*{Assumptions}
\begin{assumption}[Smoothness of $f$]\label{ass:smooth}
$f$ is convex and differentiable with $L$-Lipschitz gradient:
\[
\|\nabla f(x)-\nabla f(y)\|\le L\|x-y\|\qquad\forall x,y\in\R^{n}.
\tag{A1}
\]
Equivalently,
\[
f(y)\le f(x)+\langle\nabla f(x),y-x\rangle+\frac{L}{2}\|y-x\|^{2}.
\tag{A1'}
\]
\end{assumption}

\begin{assumption}[Convexity of $g$]\label{ass:convex}
$g$ is proper, closed, convex (possibly nonsmooth). Its proximal operator is well defined and \emph{non‑expansive}:
\[
\|\prox_{\eta g}(u)-\prox_{\eta g}(v)\|\le\|u-v\|\qquad\forall u,v\in\R^{n},\ \eta>0.
\tag{A2}
\]
\end{assumption}

\begin{assumption}[Stepsize]\label{ass:stepsize}
The stepsize satisfies
\[
0<\eta\le\frac{1}{L}.
\tag{A3}
\]
\end{assumption}

\section*{Main Convergence Theorem}
\begin{theorem}[Sublinear $O(1/k)$ convergence]\label{thm:rate}
Let $\{x_{k}\}$ be generated by \eqref{PGD} under Assumptions \ref{ass:smooth}–\ref{ass:stepsize} and let $x^{\star}\in\arg\min F$.
Then for every $k\ge1$,
\[
F(x_{k})-F(x^{\star})\;\le\;\frac{\|x_{0}-x^{\star}\|^{2}}{2\eta\,k}.
\tag{T1}
\]
Consequently,
\[
\frac{1}{k}\sum_{t=1}^{k}\bigl(F(x_{t})-F(x^{\star})\bigr)
\le\frac{\|x_{0}-x^{\star}\|^{2}}{2\eta\,k}=O\!\left(\frac{1}{k}\right).
\tag{T2}
\]
\end{theorem}

\section*{Theoretical Properties}
\begin{itemize}
\item \textbf{Stability.} The non‑expansiveness of $\prox_{\eta g}$ together with $\eta\le1/L$ guarantees that the iterates remain bounded whenever a solution exists.
\item \textbf{Hyperparameter sensitivity.} The bound \eqref{T1} improves linearly with $\eta$; however $\eta$ cannot exceed $1/L$ without losing the descent property.
\item \textbf{Comparison.} When $g\equiv0$, PGD reduces to gradient descent and the same $O(1/k)$ rate is recovered. For nonsmooth $g$, PGD attains the same rate without requiring subgradient steps.
\end{itemize}

\section*{Preliminaries (Definitions and Lemmas)}
\begin{lemma}[Optimality condition of the proximal operator]\label{lem:opt}
For any $v\in\R^{n}$ and $\eta>0$,
\[
u=\prox_{\eta g}(v)\;\Longleftrightarrow\;
\frac{1}{\eta}(v-u)\in\partial g(u).
\tag{L1}
\]
\end{lemma}
\begin{proof}
By definition of the proximal map, $u$ minimizes 
$\phi(u)=g(u)+\frac{1}{2\eta}\|u-v\|^{2}$.  
First‑order optimality yields $0\in\partial g(u)+\frac{1}{\eta}(u-v)$,
equivalently $(v-u)/\eta\in\partial g(u)$. ∎
\end{proof}

\begin{lemma}[Descent lemma for $L$‑smooth $f$]\label{lem:descent}
For any $x,y\in\R^{n}$,
\[
f(y)\le f(x)+\langle\nabla f(x),y-x\rangle+\frac{L}{2}\|y-x\|^{2}.
\tag{L2}
\]
\end{lemma}
\begin{proof}
Standard consequence of $L$‑Lipschitz continuity of $\nabla f$. ∎
\end{proof}

\section*{Main Proof (Step‑by‑Step Derivation)}
We prove Theorem \ref{thm:rate}.  
Let $x^{\star}\in\arg\min F$ and define $y_{k}=x_{k}-\eta\nabla f(x_{k})$.
Denote $s_{k+1}:=\frac{1}{\eta}(y_{k}-x_{k+1})\in\partial g(x_{k+1})$ by Lemma \ref{lem:opt}.

\begin{enumerate}
\item[Step 1.] (Optimality of the proximal step)  
From Lemma \ref{lem:opt},
\[
\frac{1}{\eta}\bigl(y_{k}-x_{k+1}\bigr)\in\partial g(x_{k+1}).
\tag{1}
\]

\item[Step 2.] (Convexity of $g$)  
For any $x$,
\[
g(x)\ge g(x_{k+1})+\bigl\langle s_{k+1},\,x-x_{k+1}\bigr\rangle.
\tag{2}
\]
Insert $s_{k+1}$ from (1):
\[
g(x)\ge g(x_{k+1})+\frac{1}{\eta}\bigl\langle y_{k}-x_{k+1},\,x-x_{k+1}\bigr\rangle.
\tag{3}
\]

\item[Step 3.] (Smoothness of $f$)  
Applying Lemma \ref{lem:descent} with $x=x_{k}$ and $y=x_{k+1}$,
\[
f(x_{k+1})\le f(x_{k})+\langle\nabla f(x_{k}),x_{k+1}-x_{k}\rangle
+\frac{L}{2}\|x_{k+1}-x_{k}\|^{2}.
\tag{4}
\]

\item[Step 4.] (Combine $f$ and $g$)  
Add $g(x_{k+1})$ to both sides of (4) and use (3) with $x=x^{\star}$:
\[
\begin{aligned}
F(x_{k+1})
&\le f(x_{k})+\langle\nabla f(x_{k}),x_{k+1}-x_{k}\rangle
+\frac{L}{2}\|x_{k+1}-x_{k}\|^{2}\\
&\qquad+g(x^{\star})-\frac{1}{\eta}\bigl\langle y_{k}-x_{k+1},\,x^{\star}-x_{k+1}\bigr\rangle .
\end{aligned}
\tag{5}
\]

\item[Step 5.] (Expand the inner product)  
Recall $y_{k}=x_{k}-\eta\nabla f(x_{k})$, thus
\[
y_{k}-x_{k+1}=x_{k}-\eta\nabla f(x_{k})-x_{k+1}
=-(x_{k+1}-x_{k})-\eta\nabla f(x_{k}).
\]
Hence
\[
\begin{aligned}
\bigl\langle y_{k}-x_{k+1},\,x^{\star}-x_{k+1}\bigr\rangle
&= -\langle x_{k+1}-x_{k},x^{\star}-x_{k+1}\rangle
-\eta\langle\nabla f(x_{k}),x^{\star}-x_{k+1}\rangle .
\end{aligned}
\tag{6}
\]

\item[Step 6.] (Insert (6) into (5))  
\[
\begin{aligned}
F(x_{k+1})-F(x^{\star})
&\le f(x_{k})-f(x^{\star})
+\langle\nabla f(x_{k}),x_{k+1}-x_{k}\rangle
+\frac{L}{2}\|x_{k+1}-x_{k}\|^{2}\\
&\qquad+\frac{1}{\eta}\langle x_{k+1}-x_{k},x^{\star}-x_{k+1}\rangle
+\langle\nabla f(x_{k}),x^{\star}-x_{k+1}\rangle .
\end{aligned}
\tag{7}
\]

\item[Step 7.] (Group the gradient terms)  
The two inner products containing $\nabla f(x_{k})$ combine:
\[
\langle\nabla f(x_{k}),x_{k+1}-x_{k}\rangle
+\langle\nabla f(x_{k}),x^{\star}-x_{k+1}\rangle
=\langle\nabla f(x_{k}),x^{\star}-x_{k}\rangle .
\tag{8}
\]

\item[Step 8.] (Use convexity of $f$)  
Convexity gives
\[
f(x^{\star})\ge f(x_{k})+\langle\nabla f(x_{k}),x^{\star}-x_{k}\rangle,
\]
or equivalently
\[
f(x_{k})-f(x^{\star})\le -\langle\nabla f(x_{k}),x^{\star}-x_{k}\rangle .
\tag{9}
\]

\item[Step 9.] (Plug (8) and (9) into (7))  
\[
\begin{aligned}
F(x_{k+1})-F(x^{\star})
&\le -\langle\nabla f(x_{k}),x^{\star}-x_{k}\rangle
+\langle\nabla f(x_{k}),x^{\star}-x_{k}\rangle
+\frac{L}{2}\|x_{k+1}-x_{k}\|^{2}\\
&\qquad+\frac{1}{\eta}\langle x_{k+1}-x_{k},x^{\star}-x_{k+1}\rangle .
\end{aligned}
\]
The first two terms cancel, yielding
\[
F(x_{k+1})-F(x^{\star})
\le \frac{L}{2}\|x_{k+1}-x_{k}\|^{2}
+\frac{1}{\eta}\langle x_{k+1}-x_{k},x^{\star}-x_{k+1}\rangle .
\tag{10}
\]

\item[Step 10.] (Rewrite the inner product)  
Using the identity $2\langle a,b\rangle =\|a\|^{2}+\|b\|^{2}-\|a-b\|^{2}$ with 
$a=x_{k+1}-x_{k}$ and $b=x^{\star}-x_{k+1}$,
\[
\begin{aligned}
2\langle x_{k+1}-x_{k},x^{\star}-x_{k+1}\rangle
&=\|x^{\star}-x_{k}\|^{2}-\|x^{\star}-x_{k+1}\|^{2}
-\|x_{k+1}-x_{k}\|^{2}.
\end{aligned}
\tag{11}
\]
Divide by $2\eta$ and insert into (10):
\[
\begin{aligned}
F(x_{k+1})-F(x^{\star})
&\le \frac{L}{2}\|x_{k+1}-x_{k}\|^{2}
+\frac{1}{2\eta}\Bigl(\|x^{\star}-x_{k}\|^{2}
-\|x^{\star}-x_{k+1}\|^{2}
-\|x_{k+1}-x_{k}\|^{2}\Bigr)\\
&= \frac{1}{2\eta}\bigl(\|x^{\star}-x_{k}\|^{2}
-\|x^{\star}-x_{k+1}\|^{2}\bigr)
+\Bigl(\frac{L}{2}-\frac{1}{2\eta}\Bigr)\|x_{k+1}-x_{k}\|^{2}.
\end{aligned}
\tag{12}
\]

\item[Step 11.] (Use stepsize condition)  
Assumption \ref{ass:stepsize} implies $\eta\le 1/L$, i.e. $\frac{L}{2}-\frac{1}{2\eta}\le0$.
Consequently the second term in (12) is non‑positive and can be dropped:
\[
F(x_{k+1})-F(x^{\star})
\le \frac{1}{2\eta}\bigl(\|x^{\star}-x_{k}\|^{2}
-\|x^{\star}-x_{k+1}\|^{2}\bigr).
\tag{13}
\]

\item[Step 12.] (Telescoping sum)  
Summing (13) for $k=0,\dots,K-1$ gives
\[
\sum_{k=0}^{K-1}\bigl(F(x_{k+1})-F(x^{\star})\bigr)
\le \frac{1}{2\eta}\bigl(\|x^{\star}-x_{0}\|^{2}
-\|x^{\star}-x_{K}\|^{2}\bigr)
\le \frac{\|x^{\star}-x_{0}\|^{2}}{2\eta}.
\tag{14}
\]

\item[Step 13.] (Deriving the $O(1/k)$ bound)  
Since the left‑hand side is a sum of $K$ non‑negative terms,
\[
K\bigl(F(x_{K})-F(x^{\star})\bigr)
\le \sum_{k=0}^{K-1}\bigl(F(x_{k+1})-F(x^{\star})\bigr)
\le \frac{\|x^{\star}-x_{0}\|^{2}}{2\eta},
\]
which yields the desired rate
\[
F(x_{K})-F(x^{\star})\le \frac{\|x^{\star}-x_{0}\|^{2}}{2\eta K}.
\tag{15}
\]
\end{enumerate}
Equation \eqref{15} is precisely the statement \eqref{T1} of Theorem \ref{thm:rate}. ∎

\section*{Rate Derivation (With Constants)}
The constant appearing in the bound is $\displaystyle C=\frac{\|x_{0}-x^{\star}\|^{2}}{2\eta}$.
If the stepsize is chosen as the maximal admissible value $\eta=1/L$, then
\[
F(x_{k})-F(x^{\star})\le \frac{L\|x_{0}-x^{\star}\|^{2}}{2k}.
\]

\section*{Special Cases}
\begin{itemize}
\item \textbf{Pure gradient descent ($g\equiv0$).} The proximal step becomes the identity, and the proof reduces to the classic $O(1/k)$ rate for smooth convex $f$.
\item \textbf{Lasso ($g(x)=\lambda\|x\|_{1}$).} The proximal operator is soft‑thresholding; the same rate holds, and the iterates automatically promote sparsity.
\item \textbf{Indicator of a convex set $C$ (projected gradient).} $g=\iota_{C}$, $\prox_{\eta g}$ is the Euclidean projection onto $C$, yielding the standard projected gradient method with identical $O(1/k)$ guarantee.
\end{itemize}

\end{document}